% =============================================================================
% Anexo A: instrumento completo das entrevistas individuais semiestruturadas
% (projeto Riachuelo, PRO3474).
%
% Documenta o roteiro na ordem em que as perguntas sao feitas ao entrevistado.
% Nao organiza as perguntas em blocos tematicos, nao discute objetivos ou
% hipoteses (isso esta na Secao 6.1, subsec:entrevistas) e nao analisa
% respostas. E apenas o instrumento.
% =============================================================================

\section{Roteiro de entrevista individual semiestruturada}
\label{anexo:roteiro_entrevistas}

Este é o Anexo A do relatório. Documenta o instrumento utilizado nas entrevistas individuais semiestruturadas descritas na Seção~\ref{subsec:entrevistas}. As perguntas aparecem na ordem em que são feitas ao entrevistado, numa sequência única, sem divisão em blocos temáticos. Perguntas de aprofundamento (\textit{probes}) aparecem em itálico, logo após a pergunta principal a que se referem, e o entrevistador as usa apenas quando fizerem sentido na conversa, não como lista fixa. Trechos entre colchetes são instruções de condução, não perguntas a fazer em voz alta. Este anexo apresenta somente o instrumento; a análise das respostas não faz parte deste documento.

\begin{enumerate}
    \item ``Para começar, me conta um pouco sobre como é a sua relação com compra de roupa. Você gosta de comprar, é mais por necessidade, funciona como um passeio?''
    \item ``Com que frequência você compra roupa em loja física, e qual unidade ou região você mais frequenta?''
    \item ``Você costuma comprar mais sozinho(a) ou acompanhado(a)?'' \textit{(Aprofundamento: ``Isso mudou nos últimos anos?'')}
    \item ``Pensa na última vez que você comprou uma peça de roupa numa loja física. Me conta essa história do início ao fim: o que te fez ir à loja, o que aconteceu lá dentro, até você sair.''
    \item ``Nesse momento, você já sabia o que estava procurando, ou foi mais explorando?''
    \item ``Como foi o momento de escolher e experimentar a peça?''
    \item ``E o pagamento, como aconteceu?'' \textit{(Aprofundamento: ``Teve algum momento que você hesitou ou quase desistiu?''; ``Alguma parte demorou mais do que você esperava?''; ``Isso é parecido com o que normalmente acontece pra você, ou foi diferente dessa vez?'')}
    \item ``Quando você está procurando alguma coisa específica numa loja, como você faz para encontrar?''
    \item ``Já aconteceu de você não encontrar o tamanho, a cor ou o modelo que queria?''
    \item ``Você pesquisa em algum canal (site, aplicativo, redes sociais) antes de ir à loja?'' \textit{(Aprofundamento: ``O que você faz quando isso acontece: pede ajuda, procura sozinho por mais tempo, desiste?''; ``Isso te faz sair da loja de mãos vazias com que frequência?''; ``Se você soubesse, antes de sair de casa, que a loja tem o que você quer, isso mudaria alguma coisa em como você compra?'')}
    \item ``Voltando à sua última compra: em algum momento um vendedor te ajudou? Como foi isso?''
    \item ``Tem alguma parte da compra em que você prefere não falar com ninguém e resolver sozinho(a)?''
    \item ``O que um vendedor resolve para você que você não resolveria sozinho?''
    \item ``Já teve uma situação em que a ajuda de um vendedor atrapalhou mais do que ajudou, ou o contrário: sentiu falta de um vendedor e não tinha ninguém por perto?'' \textit{(Aprofundamento: ``Isso muda dependendo do tipo de peça ou do valor?''; ``Isso mudou com o tempo: você busca mais ou menos ajuda hoje do que buscava antes?''; ``Se não tivesse nenhum vendedor disponível na loja inteira, o que você faria diferente?'')}
    \item ``Como foi o momento de pagar na sua última compra? Teve fila?''
    \item ``De modo geral, o que você acha do processo entre decidir que vai levar uma peça e efetivamente sair da loja com ela?''
    \item ``Já desistiu de comprar algo por causa da fila ou da demora no caixa?''
    \item ``O que acontece, na prática, entre você decidir `vou levar isso' e você sair da loja com a peça em mãos?'' \textit{(Aprofundamento: ``Isso acontece com frequência ou foi uma exceção?''; ``O que você faz quando a fila está muito grande: espera, desiste, volta depois?''; ``Alguma vez você sentiu que já tinha `decidido comprar' bem antes de conseguir efetivamente sair da loja com o produto? O que aconteceu nesse meio-tempo?'')}
    \item ``Você já usou algum tipo de autoatendimento (supermercado, farmácia, cinema) sem passar por um funcionário? Como foi?''
    \item ``Você confiaria em pagar por uma peça de roupa pelo celular, sem precisar falar com ninguém na loja?''
    \item ``Como você se sente ao sair de uma loja com uma compra? Já teve alguma experiência desconfortável com alarme, checagem de nota ou verificação na saída?'' \textit{(Aprofundamento: ``O que faria você confiar mais, ou menos, num processo assim?''; ``Isso muda dependendo do valor da peça?'')}
    \item ``Tem alguma coisa sobre comprar roupa em loja física que a gente não conversou e que você acha importante?''
    \item ``Se você pudesse mudar uma única coisa em como as lojas funcionam hoje, o que seria?''
    \item ``Já viu, em outra loja, outro país ou outro tipo de serviço, alguma coisa que resolvesse um problema parecido com os que você mencionou?''
    \item [Retomar, sem repetir de forma idêntica, a pergunta anterior sobre o que o entrevistado mudaria no processo discutido.]
    \item [Descrever o conceito de forma neutra:] ``a Riachuelo está estudando uma etiqueta de segurança na peça que passa a se comunicar com o aplicativo da loja. Ela sabe se a peça foi paga e, quando o pagamento é confirmado pelo aplicativo, ela se destrava sozinha, sem precisar passar pelo caixa.''
    \item [Perguntar a reação livre do entrevistado antes de detalhar qualquer aspecto técnico ou de segurança. Só depois, e apenas se o entrevistado não tocar no assunto espontaneamente, seguir com as perguntas de confiança e segurança abaixo.]
    \item ``O que chama sua atenção nisso, para o bem ou para o mal?''
    \item ``Em que situação você usaria isso, e em que situação não usaria?''
    \item ``Isso resolve algum dos problemas que você mencionou antes, ou resolve outra coisa?''
    \item ``O que precisaria existir, ou o que teria que te garantir, para você confiar nesse processo?''
    \item ``Você entende isso como uma mudança só na trava da peça, ou como algo que muda a experiência de compra como um todo?'' \textit{(Aprofundamento: ``Isso mudaria a forma como você usa o vendedor hoje?''; ``O que te preocuparia mais: a tecnologia falhar, ou perder contato com alguém que te ajuda quando precisa?''; ``Você imagina algum problema com uma peça `saber' que foi manuseada ou onde está dentro da loja?'')}
\end{enumerate}
